\section{A strict comparison count and simultaneous switches}
\label{sec:moment-knots}

We complete the proof of Theorem~\ref{thm:integral-schedule}.  Work under
its rank induction and suppose, for a contradiction, that no integral
schedule exists.  By Lemma~\ref{lem:hard-case}, the basic schedule has
\(l\) independent switch roots and exactly two components, denoted by
\(P,Q\), both unbalanced and of odd parity.  Every switch has positive
gap, positive total, and an interior time.  The limiting cover property
of Proposition~\ref{prop:basic-schedule} is retained.  Any realization
with fewer than \(l\) switches, or with at most \(l\) switches outside
this exceptional configuration, contradicts Lemma~\ref{lem:hard-case}.
Every feasible extra-edge link contradicts
Proposition~\ref{prop:continuation}.

Write \(\widehat P=\{\pm i:i\in P\}\) and similarly for \(Q\).  For a
signed label \(a\), put \(e_{-a}=-e_a\), \(y_a=y\cdot e_a\), and
\(v_a=(w\mu)\cdot e_a\) in the current frame.  A cross-comparison means
a pair \((a,c)\in\widehat P\times\widehat Q\).  We count each with
weight \(1/2\), so that a comparison and its signed mirror together
have weight one.  Set
\[
 r=2|P|,\qquad h=2|Q|,\qquad R=\frac{rh}{2}.
\]
Let \(p_\lambda,p_\mu,p_\xi\) be the numbers, with this same weight,
of cross ties in the initial positions, assigned moving values, and
final positions.  The middle number is independent of the time at
which the assigned values are read.

\subsection{The strict moment inequality}

Our aim is the strict inequality \(R>p_\lambda+p_\mu+p_\xi\).
The moments on the two components turn its tie counts into spectral
multiplicity counts.  For a Hermitian triple on \(\C^r\), its
\emph{profile} is the collection
\(m=(m_i(a))_{i,a}\), where \(m_i(a)\) is the multiplicity of the
numerical eigenlevel \(a\) of its \(i\)-th member.  Absent levels have
multiplicity zero.  Profiles of different sizes are compared using
the same three sets of numerical levels.

Writing \(m^P,m^Q\) for the profiles of the component moments, we
will have \(p_\lambda+p_\mu+p_\xi=(m^P\cdot m^Q)/2\).
An orbit-dimension estimate will give the required weak bound by
Cauchy--Schwarz, with a separate treatment of a moving zero level.
Without that zero, equality forces proportional profiles of full
norms \(r,h\).  We will also establish irreducibility in this equality
case, so the next lemma forces the two sizes to agree.  In the proof
of strictness, identical profiles will yield a forbidden permutation
block.

The lemma's direct-sum argument is a flag formulation of the
canonical-decomposition methods for quiver representations developed
by Kac and Schofield \cite{Kac1982,Schofield1992}. We give it directly
in terms of the three spectral filtrations.

\begin{lemma}[Profile rigidity]\label{knot:profile-rigidity}
Suppose two irreducible Hermitian zero-sum triples have proportional
profiles \(m,m'\), of sizes \(r,h\), respectively, and
\[
 \|m\|^2=r^2,\qquad \|m'\|^2=h^2.
\]
Then \(r=h\) and \(m=m'\).
\end{lemma}

\begin{proof}
\emph{Flags with scalar endomorphisms.}
For a profile \(m\), let \(\mathcal F(m)\) be the product of the three
partial-flag varieties with the prescribed upper-level dimensions.
Its complex dimension is
\begin{equation}\label{knot:flag-dimension}
 \dim_\C\mathcal F(m)
 =\frac{3r^2-\|m\|^2}{2}=r^2.
\end{equation}
The product of the corresponding Hermitian orbits has real dimension
\(2r^2\).  At an irreducible zero-sum triple, the sum map is submersive
onto the traceless Hermitian matrices: the orthogonal cokernel consists
of common commuting Hermitian matrices, which are scalar.  Its local
zero-sum fiber therefore has dimension \(r^2+1\), whereas simultaneous
unitary conjugation has orbit dimension \(r^2-1\).  Irreducibility is
open.  It follows that these fixed spectra admit arbitrarily many
pairwise unitarily inequivalent irreducible zero-sum triples.

Their upper-level flags have scalar endomorphisms, and flags from two
inequivalent triples have no nonzero morphism.  Indeed, let \(L\)
preserve all three upper-level filtrations, from a triple \((X_i)\)
to a triple \((X'_i)\).  In spectral coordinates its nonzero matrix
entries go from a level \(a\) only to levels \(a'\ge a\).  Consequently
\[
 \tr(L^*X'_iL-L^*LX_i)
   =\sum_{a,a'}(a'-a)\|L_{a',a}\|^2\ge0.
\]
The sum of these three quantities is zero.  Thus every nonzero entry
has \(a'=a\), and \(X'_iL=LX_i\) for all \(i\).  Irreducibility now
gives the assertions.  Here an invertible complex intertwiner between
irreducible Hermitian triples is a scalar multiple of a unitary
intertwiner, since \(L^*L\) commutes with the source triple.

\emph{Direct sums fill an open set.}
The Hermitian triples have supplied Hom-free spectral flags.  We now
work in the complex flag varieties: the configurations in the
direct-sum families need not themselves arise from Hermitian
zero-sum triples.  For every positive integer \(b\), direct sums of
\(b\) configurations in \(\mathcal F(m)\) having scalar endomorphisms, with arbitrary
changes of frame, contain a nonempty Zariski-open subset of
\(\mathcal F(bm)\).  To prove this, choose \(b\) pairwise Hom-free
configurations as above and examine the differential of the direct-sum
map, including the action of \(\GL(br,\C)\).  The flag tangent space
splits into ordered blocks indexed by the summands.  Its diagonal
blocks are supplied by variation of the individual flags.  Each
off-diagonal ordered block has dimension \(r^2\): the contribution
of its \(i\)-th flag is
\[
 \sum_{a<b}m_i(a)m_i(b)
 =\frac{r^2-\sum_a m_i(a)^2}{2},
\]
and summing gives \eqref{knot:flag-dimension}.  Change of
frame maps a space of \(r\times r\) matrices to that block; its kernel
is exactly the space of filtration-preserving maps between the two
summands, which is zero.  Hence this differential is surjective.
The direct-sum map is therefore dominant.  Chevalley's theorem makes
its image constructible \cite[Tag~054K]{StacksProject}; a dense
constructible subset of an irreducible variety contains a nonempty
Zariski-open subset.  Requiring scalar
endomorphisms on the summands is an open condition and includes our
chosen point, so the same assertion holds with this requirement.

We finish with the scalar-endomorphism case of the Krull--Schmidt
uniqueness argument; compare \cite[Theorem~1]{Atiyah1956}.
Choose positive integers \(b,b'\) with \(bm=b'm'\).  The two nonempty
open subsets just obtained intersect in the irreducible variety
\(\mathcal F(bm)\).  One configuration in their intersection admits
two direct-sum decompositions
\[
 V=\bigoplus_i V_i=\bigoplus_j W_j,
\]
whose summands have scalar endomorphisms, with dimensions \(r\) and
\(h\), respectively.  Inclusions and projections preserve all three
flags.  On a fixed \(V_i\), the sum of the composites
\(V_i\to W_j\to V_i\) is the identity.  One composite is therefore
a nonzero scalar.  Write it as \(gf=c\,\mathrm{id}_{V_i}\), \(c\ne0\).
The endomorphism \(fg\) of \(W_j\) is also scalar, and
\((fg)f=cf\) shows that its scalar is \(c\).  Thus \(f\) is an
isomorphism, so \(r=h\).  Proportional profiles of equal size are equal.
\end{proof}

\begin{proposition}[Strict comparison count]\label{prop:strict-count}
In the standing exceptional configuration,
\begin{equation}\label{knot:strict-count}
 C:=R-p_\lambda-p_\mu-p_\xi>0.
\end{equation}
\end{proposition}

\begin{proof}
Choose skew-symmetric zero-sum triples from the separate moment
solutions in \(P\) and \(Q\), supplied by
Lemma~\ref{lem:schedule-moment}.  Let \(m^P,m^Q\) be the profiles of
\(i\) times these triples.  Since the spectra of a real skew matrix
occur in opposite pairs,
\begin{equation}\label{knot:profile-ties}
 p_\lambda+p_\mu+p_\xi=\frac{m^P\cdot m^Q}{2}.
\end{equation}
A common nonzero skew operator commuting with either triple would give
a forbidden splitting by Lemma~\ref{lem:block-reduction}.  Thus both
common skew stabilizers vanish.

\emph{The norm bound.}
Consider one block of real dimension \(r\), write its matrices as
\(A_1,A_2,A_3\), and put \(z_i=\dim\ker A_i\).  For one matrix, a
kernel of dimension \(z_i\) contributes \(z_i(z_i-1)/2\) to its skew
stabilizer; a positive rotation level of multiplicity \(k\) contributes
\(k^2\).  The squared profile norm instead has contributions \(z_i^2\)
and \(2k^2\).  Hence twice the sum of the three stabilizer dimensions
is \(\|m\|^2-\sum_i z_i\).  The orbit sum map is submersive onto
\(\mathfrak{so}(r)\), because its orthogonal cokernel is the common
skew stabilizer.  Its kernel contains the simultaneous-conjugation
tangent, of dimension \(\dim\SO(r)\).  Comparing dimensions gives
\begin{equation}\label{knot:profile-bound}
 \|m\|^2-\sum_i z_i\le r^2-r.
\end{equation}

If the joint kernel is zero, the map
\[
 u\longmapsto(A_1u,A_2u,A_3u)
\]
is injective into the kernel of addition from
\(\bigoplus_i\im A_i\) to \(\R^r\).  This addition map is onto:
the orthogonal complement of the sum of the images is the joint
kernel.  Its kernel has dimension \(2r-\sum_i z_i\), so
\(\sum_i z_i\le r\).  Equation~\eqref{knot:profile-bound} yields
\(\|m\|\le r\).  The same proof works in odd dimension.

A joint kernel of dimension at least two would support a nonzero
commuting skew operator.  A possible joint kernel is therefore a
single line, and it must belong to the block containing the sole
possible zero of the moving weight.  Remove this line.  The reduced
triple has neither a joint kernel nor a common skew stabilizer, so
\[
 \|m\|^2
 =\|m_{\mathrm{red}}\|^2+2\sum_i z_{i,\mathrm{red}}+3
 \le(r-1)^2+2(r-1)+3=r^2+2.
\]
The moving zero-level entry in the original profile has multiplicity
two and pairs with zero in the other block's profile.  Deleting that
entry for the dot product improves the relevant squared bound to
\(r^2-2\).  Cauchy--Schwarz is then strict.  If there is a moving zero
but no joint kernel, the bounds \(\|m^P\|\le r\),
\(\|m^Q\|\le h\) are already strict in Cauchy--Schwarz, since the
moving zero entry prevents proportionality.  Thus
\(m^P\cdot m^Q<rh\) whenever the moving weight has a zero.

\emph{Excluding equality without a moving zero.}
The weak bound \(m^P\cdot m^Q\le rh\) holds by the same estimates.
Equality forces proportional profiles, \(\|m^P\|=r\),
\(\|m^Q\|=h\), and \(\sum_i z_i=r\) in the first block, with the
analogous equality in the second.  The moving matrices are invertible.
Each of the other two matrices has a kernel: if one were invertible,
the other would vanish, and the remaining pair \(A,-A\) would have
\(A\ne0\) in its common skew stabilizer.

Both triples are absolutely irreducible.  Indeed a proper common real
invariant subspace has an invariant orthogonal complement.  Each has
even dimension, since the moving restriction is invertible and skew.
Together with the other original block this gives at least three
orthogonal blocks, contrary to Lemma~\ref{lem:block-reduction}.  A
proper complex invariant subspace for the Hermitian triple gives a
nonscalar commuting Hermitian projection \(B+iD\), with \(B\) real
symmetric and \(D\) real skew.  Both commute with the real skew triple.
A nonscalar \(B\) gives the forbidden real splitting, while nonzero
\(D\) gives the forbidden skew stabilizer.

Lemma~\ref{knot:profile-rigidity} now gives \(r=h\) and identical
profiles, including their numerical levels.  Replace the second skew
triple by an orthogonal copy of the first.  Choose the isometry's
orientation to match the moving special-orthogonal orbit.  The other
two matrices impose no additional orientation restriction: each has a
kernel, and reflection in a kernel line is an orientation-reversing
commuting isometry.  Thus their orthogonal and special-orthogonal
orbits coincide.  On the sum of the two isometric copies, the skew
operator rotating one copy into the other commutes with the triple.
It gives a unitary, hence permutation-type, block in
Lemma~\ref{lem:block-reduction}, again a contradiction.  Equality is
impossible.  Equation~\eqref{knot:profile-ties} proves
\eqref{knot:strict-count}.
\end{proof}

\subsection{Core phases and the comparison function}

The moment calculation has produced the strict numerical gap
\(R-p_\lambda-p_\mu-p_\xi>0\). We now follow the cross-component
comparisons along the schedule. Strict inversions alone do not record
what happens on an interval where both positions and moving values
agree. The phases defined next provide the correction for these
double ties: such an interval will contribute \(-1\) when its two
phases are zero and \(+1\) when they are one.  The phase restriction
is a consequence of the exceptional graph, not an arbitrary choice
of a tie-breaking order.

For a signed position label \(a\), define its \emph{core phase}
\(p_a(t)\) on the open intervals between switch times.  If its axis is
off the unique cycle of its component, its phase is zero until the
switch on the first edge toward that cycle and one afterward.  If its
axis lies on the cycle, its phase is \(0,\tfrac12,1\), according as
zero, one, or two of its cycle incidences have occurred.  The phase is
unchanged on replacing \(a\) by \(-a\).

There is a useful algebraic description.  Solve the component's
root-basis equation with right-hand side \(e_a\).  The cumulative
signed \(a\)-coordinate of the edges already used is \(p_a(t)\).
For an off-cycle source, pruning the attached trees leaves a unit
contribution on its first edge toward the cycle.  At a cycle source,
the two incident cycle contributions are each one half.  This proves
the description.  Since the actual coefficients are integral on
bridges and half-odd-integral on cycle edges, phase \(1/2\) also records
that the translation term in the position formula has changed its
fractional part by a half.

\begin{lemma}[Plateaus]\label{lem:plateau}
Suppose cross labels \(a,c\) have equal position and equal assigned
moving value on an open interval between switch times.  Then
\begin{equation}\label{knot:plateau-phases}
 (p_a,p_c)=(0,0)\quad\hbox{or}\quad(1,1).
\end{equation}
\end{lemma}

\begin{proof}
The equality of moving values on distinct axes means that their
assigned shape-coordinate forms agree on the whole minimal chamber
face.  Hence their shape values agree for every time.  The initial
coordinates have the same fractional part, and equality of the current
positions implies that either both phases are half or neither is.

At any point of the plateau, reflect the entire suffix by the
reflection exchanging the two signed axes, including every later
position and switch root.  This reflection fixes the position at the
cut and fixes the current-frame shape \(w\eta(t)\) for every \(t\),
since the two coordinate forms agree.  Thus no new switch is needed
at the cut, and the reflected endpoint remains in its original Weyl
orbit.  In the root graph, cut each involved
vertex into its past and future copies, then identify opposite-time
copies crosswise.  If both phases are half, each cut opens its cycle
into a connected tree.  The crosswise identification gives a connected
\(l\)-edge graph, outside the exceptional configuration.  If neither
phase is half, one cut piece contains the cycle and the other is a
tree, possibly an isolated vertex.  The core piece uses the past copy
at phase one and the future copy at phase zero.  Opposite phases
therefore join the two cores and leave a separate tree component.
That too is excluded by Lemma~\ref{lem:hard-case}.  Only
\eqref{knot:plateau-phases} remains.
\end{proof}

For times between events, define
\begin{equation}\label{knot:J-definition}
 J(t)=\frac12\sum_{a\in\widehat P,\ c\in\widehat Q}j_{a,c}(t),
 \qquad
 j_{a,c}(t)=
 \begin{cases}
  1,&(y_a-y_c)(v_a-v_c)<0,\\
  p_a+p_c-1,&y_a=y_c\ \hbox{and}\ v_a=v_c,\\
  0,&\text{otherwise}.
 \end{cases}
\end{equation}
Use one-sided values at isolated equalities.  The direct links of
Lemma~\ref{lem:direct-links} give
\begin{equation}\label{knot:J-boundaries}
 J(0+)\ge R-p_\lambda-p_\mu,
 \qquad J(1-)\le p_\xi.
\end{equation}
Indeed strict position and velocity orders must begin discordant and
end concordant.  An initial double tie has phases zero and contributes
\(-1\), covered by subtracting both sorts of initial tie.  At the
final endpoint a double tie has phases one and contributes \(+1\);
a position tie with distinct velocities is an inversion just before
the endpoint and is also counted by \(p_\xi\).  Away from the old
switch times, any changing cross equality with distinct velocities
gives the direct crossing link.  Thus a decrease can occur only at an
old switch time, which we call a \emph{knot}.

\begin{proposition}[Monotonicity at knots]\label{prop:knot-count}
Under the standing no-integral-schedule assumption, the function
\(J\) is nondecreasing at every knot, and hence on the whole interval
in the one-sided sense used above.
\end{proposition}

We prove this by reducing each knot to disjoint transfers of two
moving values.  A loss in their strict comparison count will either
produce a local link or be compensated by a double-tie contribution.

\subsection{The form of an exact knot}

We now specify the shape genericity allowed by
Lemma~\ref{int:closedness}.  This choice concerns the original
exceptional root-basis schedule, which retains the limiting cover
property; it imposes no cover condition on later augmented schedules.
For fixed endpoint assignments and a fixed root basis, the endpoint
equation determines every coefficient.  The frames, coefficients,
and translation parts therefore range over finitely many possibilities.
Both switch conditions and equalities of signed positions consequently
give a finite collection of affine level sets for shape values.

Let \(F\) be the span of the smallest chamber face containing \(\mu\).
Its signed coordinate forms are \(\pm M_j\), for independent
variables \(M_j\), with possible repetitions.  No coordinate form
is forced to vanish: a sole zero coordinate of \(\mu\) is the free
last coordinate in type \(D\).  On distinct axes, any equality at
\(\mu\) already identifies the corresponding forms on \(F\).

Choose a polygonal shape, with slopes relatively interior in the face,
that avoids simultaneous nonconstant affine level hits having
nonproportional restrictions to \(F\), and avoids such hits at its
breaks.  Such choices are dense: after subdivision, interior break
values have open freedom in \(F\), and a generic polygonal path
avoids the finitely many affine sets of codimension at least two.

At any knot the nonzero velocity differences of its switches are
therefore multiples of one form \(\ell\).  Group signed position labels
by their assigned forms modulo \(\R\ell\).  This partition is unchanged
during the knot.  If \(\ell\) is proportional to \(M_i\pm M_j\),
\(i\ne j\), its nontrivial classes are pairs of forms and their
mirrors.  If it is proportional to \(M_i\), the only nontrivial class
is \(\{M_i,-M_i\}\).  Consequently each class has either one value or
two values \(H>L\), ordered at \(\mu\).  The sole possible exception,
\(M_i(\mu)=0\) in an opposite class, contains only the two signs of
one physical axis; it supports no root switch and may be omitted.

Across different classes, cross positions are unequal at the knot
and cross velocity values are distinct.  A swap changes their strict
inversion count only if an external spectator lies strictly between
both swapped positions and both swapped velocity values.  The
limiting cover property excludes exactly this situation by
Corollary~\ref{real:blocker}.  Thus these
comparisons are unchanged.  Classes with one value need only the
monotonicity of phases.

\begin{lemma}[Matching at a knot]\label{knot:matching}
Within each two-value class at a knot, the transfers of the high value
form disjoint matching pairs inside each component, together with
unused single labels.  In an opposite class the matching is disjoint
on physical axes, with its mirror transfers understood as the same
physical switches.
\end{lemma}

\begin{proof}
Every transfer sends \(H\) from a donor of lower position to a
recipient of higher position.  In an ordinary class, decompose this
flow into directed paths from its net donors to its net recipients.
Positions strictly increase, so no directed cycle occurs.  Shortcut
each path by its direct transfer.  The resulting matching has the
same final high and low assignments and legal positive gaps; see
Equation~\eqref{int:whole-shape}.  If the
original transfers were not already disjoint pairs, this uses fewer
switches, contrary to Lemma~\ref{lem:hard-case}.

In an opposite class each physical swap flips two axes, producing two
mirrored transfers.  Mark each axis by its initial high sign.  At any
edge occurrence, the sum of the positions in the two \emph{currently}
high signs is strictly negative.  At each vertex, pair consecutive
incidences, leaving the first unpaired if their number is odd.  The
signed-position terms in each such pair cancel, because the high sign
alternates.  The paired incidences organize the edges into paths and
cycles.  A cycle would sum strict inequalities to \(0<0\), and cannot
occur.  On a path the sum leaves just the positions of its two ends
in their initial high signs, with negative sum.  A direct flip of
these two axes is therefore feasible.  Distinct paths have disjoint
unpaired ends; their flips realize exactly the original odd-incidence
axes, hence exactly the same final sign assignment.  Unless the old
switches were already physically disjoint, this again reduces their
number.
\end{proof}

We may therefore compare units---single labels and directed pairs---of
one component with units of the other.  All local replacements below
use distinct physical axes.  Other switches at the knot are disjoint
from these axes, apart from the automatic mirrors in an opposite
class.  A transfer \(a\to b\) has root \(e_b-e_a\); all transfers in
the same two-value class have the same strictly increasing coefficient
function, namely the assigned shape value of \(H\) minus that of \(L\).

The replacements have exactly the same final assignment of \(H,L\)
as the old switches.  This preserves the whole shape, not merely its
value at the knot: the old and new final frames differ on the right
by an element of \(W_\mu\), which fixes \(F\) pointwise.  Subsequent
switches therefore have the same actual shape values and positions
at the base parameter.  The forgetting condition of
Proposition~\ref{prop:continuation} is satisfied.  We will check its
remaining hypotheses---lattice generation and feasible coefficient
order and gaps---for each extra-edge replacement.

\subsection{Pairs, single labels, and endpoint ties}

The matching lemma leaves three kinds of comparison: two unused
labels, a transfer and an unused label, and two transfers. We check
that every decrease either produces a link or is offset by the
contribution of a double tie.

Consider a transfer \(a\to b\) at positions \(x<z\), and a single
label \(c\) in the other component at position \(h\), initially high.
The low case is obtained by taking signed mirrors.  Omitting plateau terms, the strict count
decreases by one precisely when
\begin{equation}\label{knot:pair-single-overlap}
 x\le h\le z.
\end{equation}
Indeed before the knot only the comparison of \(c\) with the low
recipient can contribute; afterward only its comparison with the
new low donor can contribute.  At an endpoint equality, the distinct
approaching or departing velocities give respectively the inversion
or concordance specified by the closed inequalities in
\eqref{knot:pair-single-overlap}.

Replace the transfer by
\begin{equation}\label{knot:split}
 c\longrightarrow b,\qquad a\longrightarrow c.
\end{equation}
The final high value is at \(b\), as required.  The graph subdivides
the old edge through the other component, with the same cycle signs;
it is connected and unbalanced and generates the full root lattice.
Its base coefficients are the common old coefficient, hence
half-integral.  Let \(d_1,d_2\) be the derivatives of its two new
coefficients under a primitive half-kernel direction.  Their difference
\[
 D=d_2-d_1
\]
is nonzero: otherwise contracting the two new edges would give a
dependency of the old basis.  Choose its sign so that \(D>0\).
The common increasing coefficient function then orders the two new
times correctly.  If \(x<h<z\), both gaps are strictly positive and
\eqref{knot:split} is a forbidden feasible link.

It remains to compute the derivative of a zero gap.  For an ordered
switch root \(\beta_i\), use Equation~\eqref{int:gap}.
The two new roots are
\(\beta_1=e_b-e_c\), \(\beta_2=e_c-e_a\), with
\(\beta_1\cdot\beta_2=-1\).

\emph{A tie at the donor.}
If \(h=x\), the second new gap is zero.  Since the old root is
\(\beta_1+\beta_2\), write
\[
 d_1\beta_1+d_2\beta_2
 =d_1(\beta_1+\beta_2)+D\beta_2.
\]
The full derivative has zero root sum.  Absorbing \(d_1\) into the
old edge therefore gives an old-basis equation with right-hand side
\(-D\beta_2=D(e_a-e_c)\).  By the algebraic description of phases, its
pre-knot contributions at \(a,c\) are \(Dp_a^-,-Dp_c\).
Differentiating \eqref{int:gap}, including the first new
switch, gives
\begin{equation}\label{knot:donor-gap}
 G'_2=D(p_a^-+p_c)+d_1-d_2
      =D(p_a^-+p_c-1).
\end{equation}
Before the knot, \(a,c\) form a double plateau.  If both phases are
one, \eqref{knot:donor-gap} is positive and gives a feasible link.
Otherwise Lemma~\ref{lem:plateau} makes both phases zero.  Losing the
plateau contribution \(-1\) then compensates for the strict-count
loss of one.

\emph{A tie at the recipient.}
If \(h=z\), the first new gap is zero.  This time write the new
derivative as \(d_2(\beta_1+\beta_2)-D\beta_1\).
Absorbing \(d_2\) into the old edge gives right-hand side
\(D\beta_1=D(e_b-e_c)\) in the old-basis equation.  At the old edge,
its contribution to the \(b\)-coordinate is \(d_2\), so
\(p_b^+=p_b^-+d_2/D\).  Equation~\eqref{int:gap} gives
\begin{equation}\label{knot:recipient-gap}
 G'_1=-D(p_b^-+p_c)-d_1
      =D(1-p_b^+-p_c).
\end{equation}
The ensuing plateau of \(b,c\) has either two zero phases, giving
a feasible link, or two unit phases.  In the latter case its new
contribution \(+1\) compensates for the strict loss.  Thus every
pair--single comparison is nondecreasing when links are excluded.
Two singles can decrease only by a crossing with distinct velocities.
Their axes are unused by the knot, so this is a direct link by
Lemma~\ref{lem:direct-links}.

\subsection{Two pairs and the star replacement}

Consider \(a\to b\) at positions \(x<z\) and \(c\to d\) at
positions \(h<k\).  Before the knot the two mixed comparisons count
when \(x\le k\) and \(h\le z\); after it they count when \(x>k\)
and \(h>z\).  Thus, without plateau terms, the loss is two precisely
when the closed intervals \([x,z]\), \([h,k]\) overlap, and zero
otherwise.  The endpoint convention again follows from one-sided
times.  For example, if \(x=k\), the high donor \(a\) lies just below
the low recipient \(d\) before the knot, giving an inversion; after
the knot the same position order is concordant with their exchanged
values.

A touching mixed endpoint permits the switched matching
\[
 a\longrightarrow d,\qquad c\longrightarrow b
\]
with a zero gap and the correct other weak signs.  Deleting that
zero-gap switch leaves fewer than \(l\) edges.  Suppose henceforth
that the overlap is strict.  The two old coefficients are equal.
If they are half-odd-integral, both edges lie on their components'
cycles.  Switching the matching joins the two opened cycles into a
connected \(l\)-edge graph, again outside the exceptional
configuration.  It is immaterial whether the resulting connected
graph is balanced: a balanced graph gives the permutation reduction.

The remaining case has integral common coefficient.  Both old edges
are bridges.  For each bridge, say that its \emph{core is at an
endpoint} if the piece containing that endpoint after removal of the
bridge contains the cycle.  This switch advances the phase from zero
to one at the endpoint away from the core and leaves the core-side
endpoint phase unchanged.

Suppose the donor positions differ, say \(x<h\).  Use the star at
\(c\): send \(c\) to one of \(b,d\), then \(a\) to \(c\), then
\(c\) to the other of \(b,d\).  Strict overlap implies that all three
gaps are positive.  The new graph joins the four pieces of the two
removed bridges by this star.  It is connected, retains both
unbalanced cycles, and generates the full root lattice.  Along the
path between its two cores, the primitive half-kernel coefficients
have magnitude one; off that path the connecting-tree coefficients
vanish.  In the displayed root orientations, their possibilities are
as follows, up to reversing their common sign:
\begin{center}
\begin{tabular}{@{}ccc@{}}
\toprule
Core endpoints & \((d_{c\to b},d_{a\to c},d_{c\to d})\)
 & Wing order feasible?\\
\midrule
\(a,c\) & \((0,1,0)\) & no\\
\(a,d\) & \((0,1,1)\) & yes\\
\(b,c\) & \((1,0,0)\) & yes\\
\(b,d\) & \((1,0,-1)\) & yes\\
\bottomrule
\end{tabular}
\end{center}
To see the signs directly, only the star edges on the path from the
first core endpoint to the second are active.  Conservation at
\(c\) gives equal coefficients on \(a\to c,c\to d\), and opposite
coefficients on \(c\to b,c\to d\).  This gives exactly the table.
All three coefficients use the same increasing shape-level function.
One can therefore order the wings so that their derivatives surround
the middle derivative except when both cores are at donors.  Equal
derivatives give equal times, with the stated switch order at that
time.  In every other case this is a feasible link.

For a receiver star, if \(z<k\), use
\[
 a\longrightarrow b,\qquad b\longrightarrow d,\qquad
 c\longrightarrow b,
\]
allowing the two wing transfers into \(b\) to be interchanged.
Strict overlap makes all gaps positive.  The wing derivatives are
those of \(a\to b\) and \(c\to b\); they must bracket the middle
derivative on \(b\to d\).  If the cores are at \(b,d\), only the
middle edge lies on the path between them, so its nonzero derivative
cannot be bracketed by the two zero wing derivatives.  In every
other case, the active path uses the two wings with opposite
derivatives, the middle and one wing with equal derivatives, or one
wing alone with the other two derivatives zero.  Choosing the common
sign and the wing order makes the times weakly increasing.  Thus
only the case of two recipient cores is infeasible.  If \(k<z\),
interchange the two original pairs.

If no link exists, the donor-star and receiver-star obstructions
cannot both hold.  Thus \(x=h\) or \(z=k\), or both.  These
same-role coincidences produce double plateaus on both sides of the
knot.  If only the donors coincide, the receiver star forces both
cores to be at recipients.  Each donor phase then jumps from zero
to one, so their plateau contribution rises from \(-1\) to \(+1\).
This compensates for the strict-count loss of two.  If only the
recipients coincide, the donor star forces both cores to be at
donors, and the same compensation occurs at the recipients.  If both
pairs of endpoints coincide, the sum of the two plateau increments
is two by the bridge-phase rule, again the required compensation.  This exhausts all pair--pair
comparisons.  Mirrored transfers give the same calculations; the
factor \(1/2\) in \eqref{knot:J-definition} applies uniformly.

\begin{proof}[Proof of Proposition~\ref{prop:knot-count}]
Comparisons across distinct form classes are unchanged by the cover
property.  Within a class, Lemma~\ref{knot:matching} reduces all
comparisons to the single--single, pair--single, and pair--pair cases
just proved.  Each possible loss either gives a feasible extra-edge
link, forbidden by Proposition~\ref{prop:continuation}, or is
compensated by its plateau contribution.  The direct crossing links
exclude decreases away from knots.  Summing gives the assertion.
\end{proof}

\begin{proof}[Completion of the proof of Theorem~\ref{thm:integral-schedule}]
Proposition~\ref{prop:strict-count} and
\eqref{knot:J-boundaries} imply \(J(0+)>J(1-)\), contradicting
Proposition~\ref{prop:knot-count}.  This proves integral realizability
for the generic shapes used above.  The uniform bounded-switch
closedness of Lemma~\ref{int:closedness} passes to an arbitrary shape
with the same endpoint.  Together with the permutation case and the rank
reduction of Lemma~\ref{lem:hard-case}, this completes that theorem.
\end{proof}
